%!TEX root = ./../../springer-mono-format/book.tex
%!TEX root = ./../../thesis-trier.tex

As follows \cref{sync_max:sec:blocks},
a commutative automaton that is completely reachable
has a single state, and if it completely distinguishable
it has at most two states.\footnote{By \cref{prelim:classes:comm:prop:permSCC-regular},
	the transformation monoid of a strongly connected commutative $n$-state	automaton has size~$n$.	
	But as noted in \cref{sync_max:sec:compl_reach} the transformation monoid of a completely reachable automaton
	has size at least $2^{|Q|} - |Q|$. This also immediately implies $|Q| \le 2$ for completely reachable commutative automata.}

Next, for a fixed alphabet, we give polynomial-time algorithms to compute (i) a partially ordered automaton accepting the intersection of partially ordered commutative automata (\cref{compl:prop:intersection_waa_com}) and (ii) an automaton accepting the set of synchronizing words of a commutative automaton (\cref{prop:aut_for_sync_words}). The latter result implies that there is no exponential blow-up in the size of automata accepting this set. 
As consequences, stated in \cref{compl:cor:intersection_waa_com,cor:aut_for_sync_words},
the problems \DFAINT and \DFASWS are solvable in polynomial time for a fixed alphabet.
Lastly, in \cref{syncmax:prop:unary-sync-words} we give an expression for the set of synchronizing words of a unary automaton.

\Cref{compl:prop:intersection_waa_com,prop:aut_for_sync_words} 
were first presented in~\cite[Prop.~41 \& Lem.~43]{DBLP:journals/tcs/Hoffmann21} by the author.

\begin{proposition}
	\label{compl:prop:intersection_waa_com}
	Let $\mathcal A_i = (Q_i,\Sigma, \delta_i, q_i, F_i)$,
	$i \in \{1,\ldots,m\}$,
	be partially ordered commutative automata
	with at most $n$ states.
	Then, $\bigcap_{i=1}^m L(\mathcal A_i)$
	is recognizable by a partially ordered
	commutative automaton 
	of size $n^{|\Sigma|}$ % at most
	computable in time $O(m|\Sigma| n^{|\Sigma|+1})$.
\end{proposition} 
\begin{proof} 
	Fix an ordering $\Sigma = \{a_1, \ldots, a_k\}$ with $k \ge 1$.
	By~\cref{prelim:classes:cor:com-po} we obtain
	\begin{equation}\label{compl:eqn:com_pto_int}
		u \in \bigcap_{i=1}^m L(\mathcal A_i)
		\Leftrightarrow a_1^{\min\{n-1, |u|_{a_1}\}} \cdots a_k^{\min\{n-1, |u|_{a_k}\}} \in \bigcap_{i=1}^m L(\mathcal A_i)
	\end{equation}
	as $n \ge \max\{ |Q_i| : i \in \{1,\ldots,m\} \} - 1$. 
	Next, set $\automat{A}{Q}{\Sigma}{\delta}{q_0}{F}$
	where $Q = [n]^k$, $q_0 = (0,\ldots,0)$,
	\[
	\applytrans{\delta}{(s_1, \ldots, s_k)}{a_j} = (s_1, \ldots, s_{j-1}, \min\{n-1, s_j+1\}, s_{j+1}, \ldots, s_k),
	\]
	and 
	\[
	F = \left\{ (s_1, \ldots, s_k) \in Q \mid a_1^{s_1} \cdots a_k^{s_k} \in \bigcap_{i=1}^m L(\mathcal A_i)
	\right\}.
	\]
	Then we have
	\begin{align*}
		u \in \bigcap_{i=1}^m L(\mathcal A_i)
		& \Leftrightarrow  a_1^{\min\{n-1, |u|_{a_1}\}} \cdots a_k^{\min\{n-1, |u|_{a_k}\}} \in \bigcap_{i=1}^m L(\mathcal A_i) 
		& \tag{Eqn.~\eqref{compl:eqn:com_pto_int}}\\
		& \Leftrightarrow (\min\{n-1, |u|_{a_1}\}, \ldots, \min\{n-1, |u|_{a_k}\}) \in F 
		& \tag{Def.~$F$} \\ 
		& \Leftrightarrow \applytrans{\delta}{q_0}{a_1^{|u|_{a_1}} \cdots a_k^{|u|_{a_k}}} \in F 
		& \tag{Def.~$\delta$} \\
		& \Leftrightarrow a_1^{|u|_{a_1}}	\cdots a_k^{|u|_{a_k}} \in L(\mathcal A) 
		& \tag{Def. $L(\mathcal A)$}\\
		& \Leftrightarrow u \in L(\mathcal A). \tag{$\mathcal A$ commutative}					
	\end{align*}
	So $\bigcap_{i=1}^m L(\mathcal A_i) = L(\mathcal A)$.
		
	The automaton, without~$F$, can clearly
	be computed in time $O(n^k)$. To
	compute~$F$, 
	iterate through all $n^k$ words $a_1^{s_1} \cdots a_k^{s_k}$
	with $(s_1, \ldots, s_k) \in [n]^k$
	and test for each word whether
	$a_1^{s_1} \cdots a_k^{s_k} \in L(\mathcal A_i)$
	for all automata $\mathcal A_i$, $i \in \{1,\ldots,m\}$.
	For $i \in \{1,\ldots,m\}$, testing a single word  $a_1^{s_1} \cdots a_k^{s_k}$ of length in~$O(kn)$
	whether it is accepted by $\mathcal A_i$
	can be done in time~$O(kn)$. So testing 
	whether it is accepted by all automata $\mathcal A_i$, $i\in \{1,\ldots,m\}$,
	can be done in time $O(mkn)$. 
	In total, the algorithm runs in time $O(mkn^{k+1})$. 
\end{proof}

\begin{remark}
	The bound from \cref{compl:prop:intersection_waa_com} is tight.
	For example, set $\Sigma = \{a_1, \ldots, a_m\}$
	and for $i \in \{1,\ldots,m\}$ consider $\mathcal A_i = ([n], \Sigma, \delta_i, 0, \{n-1\})$
	with $\delta(j, a_i) = \min\{j+1,n-1\}$ and $\delta(j, a_k) = j$
	for $k \ne i$ and $j \in [n]$. We have $L(\mathcal A_i) = \{ u \in \Sigma^* : |u|_{a_i} \ge n\}$
	and $L = \bigcap_{i=1}^m L(\mathcal A_i) =  \{ u \in \Sigma^* : \forall i \in \{1,\ldots,k\} \colon |u|_{a_i} \ge n \}$.
	Let $0 \le i_1, \ldots, i_k, j_1, \ldots, j_k < n$
	where $i_r < j_r$ for some $r \in \{1,\ldots,k\}$.
	Set $u = a_1^{i_1} \cdots a_k^{i_k}$,
	$v = a_1^{j_1} \cdots a_k^{j_k}$,	
	and $w = a_1^{n - j_1} \cdots a_k^{n - j_k}$.
	Then
	\[
	uw \notin L \text{ and }
	vw \in L.
	\]
	Thus $u \not\equiv_L v$.
	This yields that the words $a_1^{i_1} \cdots a_k^{i_k}$
	for $0 \le i_1, \ldots, i_k < n$
	induce distinct Myhill-Nerode equivalence classes.
	Then every automaton accepting~$L$ has at least $n^k$ states
	and thus the bound from \cref{compl:prop:intersection_waa_com}
	is reached.	
	This also follows from results in~\cite{DBLP:journals/iandc/Hoffmann24}
	where the state complexity on commutative languages is investigated.
	The above languages were used in the proof of \cite[Thm.~1]{DBLP:journals/ijfcs/BrzozowskiL12},
	so that this proof established the above in the special case $m = 2$.
\end{remark}

\begin{corollary}	\label{compl:cor:intersection_waa_com}
	\DFAINT is solvable in time $O(m|\Sigma|n^{|\Sigma|+1})$
	for $m$ partially ordered input automata each having at most $n$ states.	
\end{corollary}
\begin{proof}
	The input automata accept a common word if and only if 
	the final state set in the automaton constructed
	in the proof of~\cref{compl:prop:intersection_waa_com}
	is non-empty.
	So, we can use the same algorithm
	as in the proof of~\cref{compl:prop:intersection_waa_com}
	and simply test whether $F \ne \emptyset$
	at the end.
\end{proof}

\begin{proposition}
	\label{prop:aut_for_sync_words}
	Let $\semiautomat{A}{Q}{\Sigma}{\delta}$
	be a commutative automaton with $n$ states.
	Then, the set of synchronizing words
	is recognizable by a partially ordered commutative automaton
	of size~$n^{|\Sigma|}$
 	computable in time $O(|\Sigma|n^{|\Sigma|+2})$.
\end{proposition}
\begin{proof}		
	By \cref{prelim:classes:comm:cor:collections-of-facts:2}, 
	for $\Syn(\mathcal A) \ne \emptyset$,
	$\mathcal A$ must have a unique trap state.
	So, first check whether $\mathcal A$
	has a unique trap state $t \in Q$, and compute it if present, 
	which can be done in time $O(n)$ by \cref{prelim:subclasses:sync:alg-unique-trap}.
	In case $\mathcal A$
	has no trap state, then $\Syn(\mathcal A) = \emptyset$
	so that a one-state automaton with empty final state set
	accepts $\Syn(\mathcal A)$.
	Otherwise,  as, by \cref{constr_sync:prop:pto_comm_rho_automata},~$\mathcal A$ is a $\rho$-automaton
	with a trap state,
	we find $\Syn(\mathcal A) = \Syn(\mathcal A / {\eqreach})$  by \cref{constr_input:prop:rho_trap_sync_factor}.
	As a quotient automaton of a commutative automaton,
	$\mathcal A / {\eqreach} = (Q / {\eqreach}, \Sigma, \delta_{\eqreach})$
	is commutative, and has the trap state $[t]_{\eqreach} = \{t\}$.
	By \cref{constr_input:prop:factor_rho_aut_pto}, 
	$\mathcal A / {\eqreach}$
	is a partially ordered automaton
	computable in linear time.	
	As, by \crefitem{prelim:thm:sync_trap_state}{prelim:subclasses:thm:trap:item:syncunique},~$[t]_{\eqreach}$ must be the unique synchronizing state, we have
	\[
	\Syn(\mathcal A / {\eqreach}) = \bigcap_{[q]_{\eqreach} \in Q / {\eqreach}} L(( \mathcal A / {\eqreach} )_{[q]_{\eqreach}, \{[t]_{\eqreach}\}}).
	\]
	By \cref{compl:prop:intersection_waa_com} 
	we can compute a partially ordered commutative
	automaton of size $n^{|\Sigma|}$ accepting this intersection involving $n$ automata,
	each having~$n$ states,	in time $O(|\Sigma|n^{|\Sigma|+2})$.
	
	The above intersection for $\Syn(\mathcal A / {\eqreach})$
	equals the same intersection but only over
	the minimal states $Q / {\eqreach}$, which might give a slightly better running time for the algorithm.
\end{proof}

We note a direct consequence of the previous result.

\begin{corollary}	\label{cor:aut_for_sync_words}
	\DFASWS\ is solvable in time $O(|\Sigma|n^{|\Sigma|+2})$
	for $n$-state commutative input automata $\semiautomat{A}{Q}{\Sigma}{\delta}$.	
\end{corollary}
\begin{proof} 
	By \cref{prop:aut_for_sync_words}, 
	we can compute an automaton of size $n^{|\Sigma|}$
	for $\Syn(\mathcal A)$ in time $O(|\Sigma|n^{|\Sigma|+2})$. Then a shortest word accepted
	by this automaton can be computed in time $O(|\Sigma|n^{|\Sigma|})$
	as a shortest path from the start state to some final state~\cite{cormen01introduction}.
\end{proof}

Martyugin~\cite[Prop.~5.2]{Mar2009b} provided a direct algorithm to decide \DFASWS  
running in time $O(|\Sigma|n^{|\Sigma|}\log n)$ for commutative input automata 
(he actually omits the factor~$|\Sigma|$ as he assumed a fixed alphabet; 
however, the proof shows that this is the running time 
when the alphabet size is taken into account).

Lastly, we note the following improvement for unary languages.

\begin{proposition} \label{syncmax:prop:unary-sync-words}
	Let $\semiautomat{A}{Q}{\{a\}}{\delta}$
	be a synchronizing unary automaton.
	Then $\Syn(\mathcal A) = a^{l} a^*$
	where $l$ is the length of a longest path traversing distinct states in~$\mathcal A$.
	In particular, $\Syn(\mathcal A) = a^{l} a^*$
	is acceptable by an automaton with at most~$l + 1\le |Q|$ states.
\end{proposition}
\begin{proof}
	A synchronizing unary automaton consists of tails all ending at a common state
	at which the single letter induces a self-loop.
	Now, it is easy to see that the set of synchronizing words
	has the form as claimed. Clearly, $l < |Q|$
	and $a^{l} a^*$ is acceptable by an automaton with $l + 1$ states.
\end{proof}

The smallest automaton whose set of synchronizing words is $a^la^*$
equals the minimal automaton accepting this language over $\{a\}$.
Thus the reset complexity and the state complexity of $a^la^*$ both equal $l+1$.
This was also stated in~\cite[Prop.~1]{DBLP:journals/corr/Maslennikova14}.